% ===== 原文 pp.64--67 =====
% 块边界：p.64 起于 6.1 节末段（上接 tr-04 到 p.63 的式 (6.5)），本块自「Correlators and conformal blocks.」
% 一段续写，式 (6.6) 为本块第一条编号公式，与 (6.5) 相接；本块不重写 §6、6.1 的标题。

% [原文 p.64]
\para{关联函数与共形块。}
主算子的 $n\ge 4$ 点函数于是被共形不变性固定 [253]。为具体起见，考虑一个 4 点函数，
并插入一套完备态，再按 $\mathrm{Vir}\times\overline{\mathrm{Vir}}$ 的表示把它们分门别类：
\begin{equation}
\langle V_1V_2V_3V_4\rangle
=\sum_{\substack{V,V'\ \text{primaries}\\ Y,\bar Y,Y',\bar Y'\ \text{Young diagrams}}}
\langle V_1V_2\,\mathbb{L}_{-Y}\mathbb{L}_{-\bar Y}|V\rangle\;
g^{-1}(Y,\bar Y,V;Y',\bar Y',V')\;
\langle V'|\,\mathbb{L}_{Y'}\mathbb{L}_{\bar Y'}V_3V_4\rangle,
\tag{6.6}
\end{equation}
其中 $g^{-1}(Y,\bar Y,V;Y',\bar Y',V')$ 表示以
$\langle V'|\mathbb{L}_{Y'}\mathbb{L}_{\bar Y'}\mathbb{L}_{-Y}\mathbb{L}_{-\bar Y}|V\rangle$
为分量的（矩阵）之逆的分量。把全部 Virasoro 生成元都翻译成作用在式 (6.5) 中位置依赖性上的微分算子，即得
\begin{equation}
\langle V_1V_2V_3V_4\rangle
=\sum_{V_5,V_6\ \text{primaries}}
C_{125}\,g^{-1}_{56}\,C_{634}\,
\mathcal{F}(h_1,\dots,h_5;z_1,\dots,z_4)\,
\overline{\mathcal{F}(\bar h_1,\dots,\bar h_5;\bar z_1,\dots,\bar z_4)}.
\tag{6.7}
\end{equation}
相差一些无关紧要的因子，这个（局部）全纯因子 $\mathcal{F}$ 称为\emph{共形块}（conformal block）。
它完全由外算子（external operators）的维数 $h_1,\dots,h_4$，
以及作为完备态的一部分插入的内算子（internal operator）的维数 $h_5=h_6$ 决定。

我们本可以把完备态插到别处，即换一种方式把各算子 $V_i$ 分到插入态的两侧。
更一般地，把一个球面 $n$ 点函数归结为常数 $g$、$C$ 与共形块的所有可能方式，
对应于把 $n$ 穿刺球面分解成若干三穿刺球面的所有可能方式：
每个把球面切成碎块的闭合圈上都插入一套完备态。
无论哪种情形，共形块都是纯表示论的对象；
它们依赖于 $n$ 个外算子的维数，以及沿割线插入的 $n-3$ 个内算子的维数。

% [Tutor Note: p.64 式 (6.8) 的量纲核对：作用量 $S=\frac{1}{4\pi}\int\dd^2z\sqrt{g}(\partial_\nu\phi\partial^\nu\phi+QR\phi+4\pi\mu\ee^{2b\phi})$ 中
%   二维标量按通常归一化无量纲，故 $b$、$Q$ 无量纲，而 $\ee^{2b\phi}$ 前的 $\mu$ 必须带 $(长度)^{-2}$（与 $\partial\phi\partial\phi$、$R\phi$ 两项同量纲），
%   这一点也是 (6.9) 里 $\lambda=(\pi\Gamma(b^2)\mu/\Gamma(1-b^2))^{1/b}$ 带量纲、而 $\lambda^{Q-2\alpha}$ 的幂次要靠 $b$ 与 $Q=b+1/b$ 的约定才配平的由来。
%   另核 $c=1+6Q^2\ge 25$：由 $Q=b+1/b\ge 2$（$b>0$）得 $c\ge 1+24=25$，等号在 $b=1$，正是 Liouville 的中心荷下界。]

\para{Liouville 理论。}
Liouville 理论描述单个标量场，其作用量为
\begin{equation}
S[\phi]=\frac{1}{4\pi}\int \dd^2z\,\sqrt{g}\,
\big(\partial_\nu\phi\,\partial^\nu\phi+QR\phi+4\pi\mu\,\ee^{2b\phi}\big)
\tag{6.8}
\end{equation}
其中 $R$ 是 Ricci 标量。只要 $Q=b+1/b$，该理论就是共形的，
其全纯能动张量为 $T=(\partial\phi)^2+Q\partial^2\phi$，中心荷为 $c=\bar c=1+6Q^2\ge 25$。
虽然宇宙学常数 $\mu$ 看上去像个耦合常数，但事实表明它在关联函数中只以平凡的方式出现：
真正有意思的是对 $b>0$ 的依赖，其中 $b\to 0$ 是半经典极限。
Liouville 共形场论具有一个（非显然的）对偶 $b\to 1/b$，
同时保持 $\lambda=\big(\frac{\pi\Gamma(b^2)}{\Gamma(1-b^2)}\mu\big)^{1/b}$ 不变。
可以验证 $V_\alpha=\,:\!\ee^{2\alpha\phi}\!:\,$ 是左/右行维数为
$h(\alpha)=\alpha(Q-\alpha)=Q^2/4+P^2$ 的共形主算子，其中 $\alpha=(b+1/b)/2+iP$、$P\in\R$。
不变性 $h(Q-\alpha)=h(\alpha)$ 提示作如下认同：$V_\alpha=R(\alpha)V_{Q-\alpha}$。
反射系数可以（用共形自举）确定为
\begin{equation}
R(\alpha)=-\lambda^{Q-2\alpha}\,
\frac{\Gamma\big(b(2\alpha-Q)\big)\Gamma\big(\tfrac{1}{b}(2\alpha-Q)\big)}
{\Gamma\big(b(Q-2\alpha)\big)\Gamma\big(\tfrac{1}{b}(Q-2\alpha)\big)}.
\tag{6.9}
\end{equation}

% [原文 p.65]
两点函数于是为
\begin{equation}
\langle V_{\alpha_1}V_{\alpha_2}\rangle
=\delta_{\alpha_1+\alpha_2=Q}+R(\alpha_1)\,\delta_{\alpha_1=\alpha_2}.
\tag{6.10}
\end{equation}
三点函数已知由 Dorn--Otto--Zamolodchikov--Zamolodchikov（DOZZ）公式给出 [254, 255]
\begin{equation}
\begin{aligned}
C_{\alpha_1\alpha_2\alpha_3}
&=\langle V_{\alpha_1}V_{\alpha_2}V_{\alpha_3}\rangle\\
&=\frac{(b^{2/b-2b}\lambda)^{Q-\alpha_1-\alpha_2-\alpha_3}\,
\Upsilon'_b(0)\,\Upsilon_b(2\alpha_1)\Upsilon_b(2\alpha_2)\Upsilon_b(2\alpha_3)}
{\Upsilon_b(\alpha_1+\alpha_2+\alpha_3-Q)\Upsilon_b(\alpha_1+\alpha_2-\alpha_3)
\Upsilon_b(\alpha_2+\alpha_3-\alpha_1)\Upsilon_b(\alpha_3+\alpha_1-\alpha_2)}.
\end{aligned}
\tag{6.11}
\end{equation}
四点函数例如读作
\begin{equation}
\langle V_{\alpha_1}(0)V_{\alpha_2}(q)V_{\alpha_3}(1)V_{\alpha_4}(\infty)\rangle
=\frac{1}{2}\int_{Q/2+i\R}\dd\alpha_s\;
C_{\alpha_1\alpha_2\alpha_s}C_{(Q-\alpha_s)\alpha_3\alpha_4}
\Big|q^{h_s-h_1-h_2}\big(1+O(q)\big)\Big|^2
\tag{6.12}
\end{equation}
其中因子 $1/2$ 抵消认同 $\alpha_s\sim Q-\alpha_s$ 带来的重复计数，
而 $1+O(q)$ 表示 $q$ 的正整数次幂的无穷级数，即归一共形块。

% [Tutor Note: p.65 式 (6.11) 与 (6.14) 的一致性核对方法：按附录 A，$\Upsilon_b(x)=1/(\Gamma_b(x)\Gamma_b(b+b^{-1}-x))$，故 $\Upsilon_b(x)=\Upsilon_b(Q-x)$，这正是正文所用的不变性。把 (6.14) 的 $\prod_{\pm\pm}$ 按四种符号展开，宗量依次为
%   $\alpha_1+\alpha_2+\alpha_3-Q$、$\alpha_1+\alpha_2-\alpha_3$、$\alpha_1-\alpha_2+\alpha_3$、$\alpha_1-\alpha_2-\alpha_3+Q$。
%   前三个直接就是 (6.11) 分母的对应因子；只有第四个需用一次反射 $\Upsilon_b(Q-x)=\Upsilon_b(x)$ 化成 $\Upsilon_b(\alpha_2+\alpha_3-\alpha_1)$，即第四个因子。
%   （易错：不能把 $\alpha_1\pm(\alpha_2-Q/2)\pm(\alpha_3-Q/2)$ 两两“由反射性互相等同”——它们相差 $2(\alpha_1-Q/2)$，并不相等；只有 $x\leftrightarrow Q-x$ 的单次反射才合法。）
%   再由 (6.15) 的归一化把 $\Upsilon'_b(0)\Upsilon_b(2\alpha_1)\Upsilon_b(2\alpha_2)\Upsilon_b(2\alpha_3)$ 吸收干净，这就是 AGT 对应的第一处精确匹配。]

\begin{exercise}{6.3}
1. 检验 DOZZ 公式 (6.11) 尊重所期望的 $b\to 1/b$ 对偶，以及对任意 $i$ 的对称性
$\alpha_i\to Q-\alpha_i$，至多相差相应的反射系数。

2. 利用附录 A 中列出的 $\Upsilon_b$ 的性质，证明在固定的普遍 $\alpha_1,\alpha_2$ 下，
$C_{\alpha_1\alpha_2\alpha_3}$ 的 $\alpha_3\to 0$ 极限为零；证明当 $\alpha_1=\alpha_2$ 时该极限为无穷大，
而 $(\alpha_3/2)C_{\alpha\alpha\alpha_3}\to g_{\alpha\alpha}=R(\alpha)$。
\end{exercise}

\subsection{寻找 AGT 字典}

我们期待存在如下形式的关系
\begin{equation}
Z_{S^4_b}\big(\mathrm{T}\big(\fsu(2),C,m\big)\big)
=\big\langle \widehat{V}_{\alpha_1}(z_1)\dots\widehat{V}_{\alpha_n}(z_n)\big\rangle_{\overline{C}}
\tag{6.13}
\end{equation}
对任意穿刺点数 $n$ 成立，其中 $\widehat{V}_{\alpha_i}(z_i)$ 是 6d 理论的余维 2 算子约化到点上的结果。
本节中我们利用已知的 $S^4_b$ 配分函数来确定：相应的 2d CFT 就是上面描述的 Liouville 共形场论，
而 $\widehat{V}_\alpha$ 是顶点算子 $V_\alpha$ 的适当重标度。

\para{三点函数与归一化。}
2d CFT 由其谱（主算子的左、右共形维数）与 OPE 结构常数（等价地，共形主算子的三点函数）刻画。
在用 $\cX(\fsu(2))$ 构造 class S 理论时，与一个穿刺点相关联的数据是一个质量参数 $m\in i\R/\Z_2$。
因此我们想要的是带连续参数的局域算子 $V$。
为与先前的记号一致，我们把这个（无量纲的）参数记作 $\alpha=Q/2+rm$，其中 $Q=b+1/b$。

% [原文 p.66]
$\widehat{V}_\alpha$ 的共形维数暂且放一放；先看三点函数。
我们知道，与三穿刺球面相关联的理论是一个三基本表示半超多重态。
其配分函数是超多重态的单圈行列式 (5.25)，于是三点函数为
\begin{equation}
\begin{aligned}
\widehat{C}_{\alpha_1\alpha_2\alpha_3}
&:=\langle \widehat{V}_{\alpha_1}\widehat{V}_{\alpha_2}\widehat{V}_{\alpha_3}\rangle
=\prod_{\pm\pm}\frac{1}{\Upsilon_b\big(\alpha_1\pm(\alpha_2-Q/2)\pm(\alpha_3-Q/2)\big)}\\
&=\frac{1}{\Upsilon_b(\alpha_2+\alpha_3-\alpha_1)\Upsilon_b(\alpha_3+\alpha_1-\alpha_2)
\Upsilon_b(\alpha_1+\alpha_2-\alpha_3)\Upsilon_b(\alpha_1+\alpha_2+\alpha_3-Q)},
\end{aligned}
\tag{6.14}
\end{equation}
其中我们用到了不变性 $\Upsilon_b(x)=\Upsilon_b(Q-x)$。
这与 DOZZ 公式 (6.11) 的分母精确吻合，
而分子（除一个与 $\alpha$ 无关的因子外）可以被下述归一化吸收：
\begin{equation}
\widehat{V}_\alpha
=\frac{(b^{2/b-2b}\lambda)^{\alpha-Q/2}}{\Upsilon_b(2\alpha)}\,V_\alpha.
\tag{6.15}
\end{equation}
在这一归一化下可以验证 $\widehat{V}_\alpha=\widehat{V}_{Q-\alpha}$，且两点函数读作
\begin{equation}
\widehat{g}_{\alpha\alpha'}
=\langle\widehat{V}_\alpha\widehat{V}_{\alpha'}\rangle
=\frac{\delta_{\alpha+\alpha'=Q}+\delta_{\alpha=\alpha'}}
{\Upsilon_b(Q-2\alpha)\Upsilon_b(2\alpha-Q)}.
\tag{6.16}
\end{equation}

\para{四点函数与维数。}
为确定 $\widehat{V}_\alpha$ 的共形维数，我们考虑一个四穿刺球面，
并沿一条适合 $q\to 0$ 极限的通道把它切开，其中 $q$ 是四个穿刺点的交比。
与四穿刺球面对应的规范理论是 $\fsu(2)$、$N_f=4$ 的 SQCD，
其配分函数用超对称局域化算得，形式为 (5.28)
\begin{equation}
Z_{S^4_b}=\int_{i\R/\Z_2}\dd a\,|q|^{r^2|a|^2}
Z_{\text{one-loop}}(a)Z_{\text{inst}}(a,q)\overline{Z_{\text{inst}}(a,\bar q)}.
\tag{6.17}
\end{equation}
在 $q\to 0$ 极限下 $Z_{\text{inst}}\to 1$。
这一表达式应当与 2d CFT 中四点函数的分解作比较：
\begin{equation}
\begin{aligned}
&\langle\widehat{V}_{\alpha_1}(0)\widehat{V}_{\alpha_2}(q)
\widehat{V}_{\alpha_3}(1)\widehat{V}_{\alpha_4}(\infty)\rangle\\
&\qquad=\int_{Q/2+i\R/\Z_2}\dd\alpha\;
q^{h(\alpha)}\bar q^{\bar h(\alpha)}
\frac{\widehat{C}_{\alpha_1\alpha_2\alpha}\widehat{C}_{\alpha\alpha_3\alpha_4}}{\widehat{g}_{\alpha\alpha}}
\,\mathcal{F}(\alpha_i,\alpha;q)\overline{\mathcal{F}(\alpha_i,\alpha;\bar q)}
\end{aligned}
\tag{6.18}
\end{equation}
其中 $\mathcal{F}$ 是共形块，它们（反）全纯地依赖于交比 $q$，并在 $q\to 0$ 时趋于 1。

我们已经把三点函数 $C$ 认同为超多重态的单圈行列式。
反过来，逆两点函数 $\widehat{g}^{-1}_{\alpha\alpha}$ 等于矢量多重态的单圈行列式。
因此很自然地期待共形块与瞬子配分函数相符，
并把 $q$ 的幂次------即 $h(\alpha)=\bar h(\alpha)=Q^2/4+P^2$ 与 $r^2|a|^2$------
在相差一个无关紧要的 $Q^2/4$ 平移的意义下等同起来。

% [Tutor Note: p.66 对 (6.12) 与 (6.18) 的积分域要留意：前者是 $Q/2+i\R$，后者是 $Q/2+i\R/\Z_2$。
%   差别来自 (6.16) 的 $\widehat g_{\alpha\alpha'}$ 已把 $\alpha\sim Q-\alpha$ 的反射包含在内
%   （其分母在 $\alpha\to Q-\alpha$ 下不变），故 (6.18) 不再需要 (6.12) 里那个抵消重复计数的 $1/2$。
%   核对方法：把 (6.16) 代入 (6.18) 的被积函数，作 $\alpha\to Q-\alpha$ 可见被积式不变，两式相差的正是这个 $\Z_2$。]

% [原文 p.67]
\para{共形块与 AGT 的证明。}
要检验 AGT 字典，剩下的关键一环是验证共形块确实与瞬子配分函数
\footnote{有趣的是，$\cN=4$ SYM 的情形在文献中引起过一些早期的混乱，这一点在 [256] 中被澄清：
该理论有 $Z_{\text{inst}}=1$，它对应于一个带非平凡顶点算子插入 $\alpha=Q/2$ 的环面。}
相符，而这一点已在低阶（按 $q$ 的幂次）被检验过 [5, 93, 257--263]。
已有许多途径（例如见 [124, 第 5.3 节] 的简述）。
一类途径依赖于在瞬子模空间上展示 Virasoro 代数（以及许多推广）的作用。
早期例子见 [264, 265]，推广见 [14, 29, 82, 266--284]。
特别地，可以构造 [14, 29, 245, 285, 286]（另见 [287--289]）
$|\widehat{V}_\alpha\rangle$ 的共形状射态的一组正交归一基，
使得按 (6.6) 把这些态插入四点函数时，逐项给出 Nekrasov 瞬子配分函数
按 U(1)-不变的类点瞬子构型求和的表达式。
Virasoro 代数与 W 代数也在 6d 语境中出现 [290--293]。
另见我们在第 101 页第 10 节对更精巧对称代数的讨论。

递推关系在 [227, 271, 294--296] 中被研究。
一个困难例如是瞬子展开中出现的伪极点，它们在把全部贡献求和后消失 [297]。
大 $c$ 极限在 [298, 299] 中研究。
基于 Dotsenko--Fateev 的 CFT 关联函数表示的自由场途径见 [289, 300--306]。
AGT 字典的弦论推导（来自一个 5d 推广）见 [67, 307]，综述见 [308]。
一条相当不同的途径是把共形块与瞬子配分函数都刻画为 Riemann--Hilbert 问题的解 [309]。

\subsection{从 6d 得到 Liouville}

我们已经论证，AGT 对应所涉及的 2d 理论是 Liouville 共形场论，
并且 AGT 对应的众多检验都证实了这一点。我们能直接从 6d 看到它吗？

\para{途径。}
$C$ 的度规中保持其共形类（等价地，保持其复结构）的形变，
关于我们用于超对称局域化的超荷 $Q$ 是 $Q$-恰当的 [18]。
因此这类形变不影响 6d 理论的配分函数，
而后者可以在 $C$ 比 $S^4_b$ 无穷小或无穷大的极限下计算。
重要的是，即便存在诸如我们在第 8 节考虑的圈、面或壁这类 $Q$-闭的可观测量，这一论证依然成立。
记住 6d 理论是共形不变的，这些极限就等价于沿其中任一因子作维数约化。
于是我们应当期待：沿 $S^4_b$ 对 6d 理论 $\cX(\fsu(2))$（或 $\cX(\fsu(N))$）作维数约化，
就得到 Liouville 共形场论（或其高阶推广，即 7.1 小节进一步讨论的 Toda 共形场论）。

在 [18, 310] 中，Córdova 与 Jafferis 分三步完成了这一约化：
